\chapter{Come funziona texweb}\label{ch:intro}
Questi appunti si scrivono in \LaTeX{}, come sempre: lo stesso sorgente produce il PDF con \texttt{pdflatex} e il sito con \texttt{texweb dev}. Riferimenti, citazioni e numerazione funzionano in entrambi, ad esempio il capitolo~\ref{ch:conv} parla di convoluzione e cita \cite{Goodfellow-et-al-2016}.

\section{Cosa si può scrivere}
Il dialetto contiene i comandi di tutti i giorni: \textbf{grassetto}, \emph{corsivo}, \texttt{codice}, note a piè di pagina\footnote{Come questa, numerate per pagina.} e formule come $f(x) = \max\{0, x\}$ oppure
\begin{equation}\label{eq:loss}
  \mathcal{L}(\vartheta) = \frac{1}{N}\sum_{i=1}^{N} \big(y^{(i)} - h(x^{(i)}; \vartheta)\big)^2, \qquad \vartheta \in \R^d.
\end{equation}
L'equazione~\eqref{eq:loss} è la MSE.

\begin{definizione}[Funzione di loss]\label{def:loss}
Una \textbf{funzione di loss} misura l'errore commesso dal modello nel predire $y$ a partire da $x$.
\end{definizione}

La definizione~\ref{def:loss} è numerata per capitolo, come nel PDF.

\begin{itemize}
  \item gli elenchi puntati,
  \item gli elenchi numerati e le tabelle,
  \item le figure TikZ, compilate da \LaTeX{} in SVG.
\end{itemize}

\begin{table}[htbp]
\centering
\begin{tabular}{@{}lrr@{}}
\toprule
Modello & Parametri & Layer \\
\midrule
MLP & 397\,510 & 2 \\
CNN & 156 & 1 \\
\bottomrule
\end{tabular}
\caption{Parametri del primo layer per MNIST.}
\label{tab:param}
\end{table}

Nella tabella~\ref{tab:param} si vede il risparmio di parametri.
